\section{What did catch it}
\label{sec:caught}

The errors were found on a second pass over the same data, undertaken because
the results felt thin rather than because anything had flagged them. Three
checks did the work, and none compares two derivations of a quantity. Each goes
back to the source's own account of itself.

\paragraph{Read the field documentation, not the data.}
The vocabulary had been inferred from the values that appeared in the data,
which is how docking and attachment came to be read as statements about where an
object is. GCAT publishes a page defining every status value. It says that a
record represents a phase in an object's flight history and that the status
names the event that ends the phase. Reading that separated ``the object is
somewhere'' from ``this record stops here.''

\paragraph{Decompose every unexplained residue.}
The first write-up reported 769 of its disagreements as unexplained. Stating a
residue feels like rigour, which is why errors hide there comfortably. Asking what
the 769 consisted of ended most of the finding: 640 paired a CelesTrak impact or
landing with a GCAT transition code, 500 of them docking or attachment. Gemini~8
is the clearest case. GCAT ends the record at docking with the Agena target
vehicle and CelesTrak records the landing, and both are right about different
events in one mission.

\paragraph{Ask whether the source already models what you accuse it of omitting.}
The tracking-loss finding rested on an implicit claim that CelesTrak has no way
to record lost tracking. Checking that claim directly, rather than the count
derived from it, found the field in use.

Redundancy asks whether two routes from $A$ to $B$ agree. These three ask
whether $A$ is what we think it is.

\section{The correction was wrong too}
\label{sec:second}

The first correction moved \texttt{C} (collided) into the destroyed set, yet
GCAT's definitions allow an object to survive a collision. Checking that one
entry against the source found more than the entry itself.

\paragraph{The source's history of an object is in a file we never read.}
GCAT's definitions say a collision destroys the object ``in which case there is
no subsequent phase'', while a survivor starts a new phase; and that after
\texttt{E} (exploded) the ``next phase of this object is a debris fragment.''
So \texttt{E} marks a boundary, not a fate, and \texttt{C} is a fate only when
nothing follows it. Whether anything follows cannot be seen in the files the
pipeline read, because each holds exactly one record per object. GCAT's own
index says where the rest is: an event catalogue described as holding ``later
phases in history for objects in primary catalogs'', listed on the same page as
the four files we fetched, under a different heading.

\paragraph{Checking the correction against those histories.}
We appended each object's event-catalogue phases to its record, took the latest
phase as GCAT's current account of the object, resolved docked and attached
objects through the object they joined, and recomputed the comparison
(Appendix~\ref{app:history}). Of the 23 objects coded \texttt{C} that both
catalogues carry, 21 have a later phase; of the 78 coded \texttt{E}, 66 do. The
count falls from 261 to 220. It loses 42 objects, all 30 explosions and 12 of the
14 collisions in the 261. For 40 of them GCAT records a later phase that is still
in orbit, which is where CelesTrak has them; among these are Cosmos~2251 and
Iridium~33 after their 2009 collision and Fengyun-1C after the 2007
anti-satellite test. The other two are explosions with no later phase recorded,
on which GCAT makes no present claim. One object enters, a payload whose later
phase ends in a deorbit that CelesTrak does not record.

A second reading from the source agrees. GCAT publishes a derived catalogue of
the most recent phase of every object. Read against it, the count is 219, and 208
objects are common to both readings. Every one of the 23 differences has a named
cause: 12 are reentries recorded by one register and not yet by the other in the
six weeks between retrievals, 3 are GCAT record revisions in that window, and 8
are deep-space objects the pipeline sets aside by design and the derived
catalogue follows. Both readings drop all 42 objects. We treat 220 as the
reference count, with the caveat that its event-catalogue input was retrieved on
29~September (GCAT's data update of 28~September) rather than 18~August.

\paragraph{What this shows.}
The first correction was made by reading the documentation, as the paper
recommends, and it still misread two codes and missed the file that would have
settled them. It was a new reading of the vocabulary, and the pipeline's gate
and our three checks both inherited it. None of them
flagged the 42 objects (Section~\ref{sec:gates}); a later line-by-line reading
of the source's definitions did.
